\begin{problem}{\bf CBC encryption} (5 points)\\
Recall that AES is a block cipher with plaintext and ciphertext blocks
of length 128 bits, and key space $\{0,1\}^{128}$. Plaintext messages
$m$ that are longer than 128 bits are broken into blocks:
$m = (m_1,m_2,\ldots,m_t)$ where each block $m_i$ is 128 bits long.
In the CBC mode of operation, $m$ is encrypted by first selecting
$c_0 \in_R \{0,1\}^{128}$ and then computing
$c_i = \mbox{AES}_k(m_i \oplus c_{i-1})$ for $1 \leq i \leq t$;
the ciphertext is $c=(c_0,c_1,c_2,\ldots,c_t)$.
(Note that a new $c_0$ is selected each time a message is encrypted.)\\
Show that CBC encryption is \emph{not} semantically secure against
a \emph{chosen-ciphertext attack}.
\end{problem}

\begin{loigiai}
\end{loigiai}
